\documentclass[10pt,a4paper]{amsart}
\usepackage[margin=19mm]{geometry}
\usepackage{amsmath,amssymb,mathtools}
\usepackage[scr=rsfs,cal=euler]{mathalfa}
\usepackage{xfrac}
\usepackage[unicode,hidelinks,bookmarks=false]{hyperref}
\newcommand{\ie}{i.e.\ }
\newcommand{\cf}{cf.\ }
\newcommand{\resp}{respectively\ }
\newcommand{\Subsection}{\subsection}
\providecommand{\nobreakdash}{\nobreak\hbox{-}}
\setlength{\emergencystretch}{2em}
\multlinegap0pt
\usepackage{bm}
\usepackage{amssymb}
\usepackage{tikz}
\usepackage{float}
\usepackage{braket}
\newtheorem{theorem}{Theorem}
\title{Floquet Codes from Derived Semi-Regular Hyperbolic Tessellations \\ on Orientable and Non-Orientable Surfaces}
\author{Douglas F. Copatti, Giuliano G. La Guardia, Waldir S. Soares, Edson D. Carvalho, Eduardo B. Silva}
\date{Source: arXiv:2603.29811v1 (CC BY 4.0)}
\begin{document}
\maketitle

\noindent\emph{Source and licence.} Floquet Codes from Derived Semi-Regular Hyperbolic Tessellations \\ on Orientable and Non-Orientable Surfaces. Version arXiv:2603.29811v1 (CC BY 4.0), distributed by arXiv under the Creative Commons Attribution 4.0 International licence, http://creativecommons.org/licenses/by/4.0/, as recorded in the arXiv licence metadata for this exact version and shown on the abstract page https://arxiv.org/abs/2603.29811v1. Full licence notice: \texttt{licenses/arxiv-2603.29811-CC-BY-4.0.txt}.\par\smallskip

\noindent\textit{Editorial scope.} Counted unit: the article's theorem or\_same\_non\_or\_qcc (source0.tex lines 1581--1583). The article's complete original proof of it is delivered with it (source0.tex lines 1585--1595, one \texttt{proof} environment of 11 lines). No mathematics is written, repaired or re-derived: the statement and the proof are the article's own lines, in the article's own notation, and no retained line is substituted or re-flowed.  No further source-local context is needed: every cross-reference of every retained line resolves inside the two delivered spans.  Nothing was omitted from any retained span, so the package carries no omission record.  No retained line contains a citation, so the wrapper carries no bibliography and no bibliography entry.  No retained line contains a figure environment, an image-inclusion command or drawing code, so no image is delivered and the package declares no asset dependency.  Editorial formatting only, in addition: an AMS \texttt{amsart} preamble that restates the article's own declarations and macros used by the retained lines, namely the environments \texttt{theorem} on separate counters, together with the standard packages those lines need.  The retained environments print in this document as Theorem 1 in order of appearance, whereas the article prints the same items with the numbering of its own full document.  The complete native source of the article as distributed in the arXiv e-print for this exact version is bundled unchanged as \texttt{sources/197542/source0.tex}.
\begin{theorem}\label{or_same_non_or_qcc}
	Let $S_{g}$ and $S_{h}$ be surfaces with non-orientable and orientable genus $g$ and $ h $, respectively. If $\chi(S_{g}) = \chi(S_{h})$, then the Floquet codes constructed are the same for both surfaces. 
\end{theorem}

\begin{proof}
	We know that the fundamental region of $S_{h}$ is a $4h$-gon polygon. Similarly, the fundamental region of $S_g$ is a $2g$-gon polygon. If $\chi(S_{g}) = \chi(S_{h})$, then $g=2h$. Hence, the fundamental region for both surfaces, $S_g$ and $S_h$ are the same $2g$-gonal polygon with different orientations on the boundary edges. Therefore, both the surfaces admit the same existing $[p,q,r]$ tessellation, and hence, the stabilizer operators associated with the tessellation $[p,q,r]$ for both the surfaces are the same since both are independent of the orientation of the boundary edges. Therefore, the code space and the encoding space associated with the tessellation $[p,q,r]$ on both surfaces are the same.
	
	We have that the code parameter $k$ is given by $k=4-2\chi$, therefore, $k$ is the same for both surfaces $S_g$ and $S_{h}$ if $\chi(S_g)=\chi(S_{h})$.
	
	We see that $l(p,q)$, $D(p,q)$ and $AR(p,q)$ are not depending on the genus of the surface. Therefore, if $\chi(S_g)= \chi(S_{h})$, then $l(p,q)$, $D(p,q)$ and $AR(p,q)$ are same for both surfaces $S_g$ and $S_{h}$. Now, $\chi(S_g)=  \chi(S_{h})$ implies $g=2h$, and this leads to
	\[
	d_h=2arccosh\Bigg[\frac{\cos{(\frac{\pi}{4h})}}{\sin{(\frac{\pi}{4h})}}\bigg].
	\]
	This concludes that $d_{min}$ is the same for the surface $S_{h}$.
\end{proof}

\end{document}
